Hospital management is a highly significant domain because it brings together patients, medical specialists, hospital staff and clinical data. Without proper software and databases, paper-based management leads to conflicts between people, a poor patient experience and inefficient operations. This report presents \textbf{Care \& Cure HMS}, a full-stack Hospital Management System built with React.js, TypeScript, Node.js, Express.js and MySQL, and describes how the first version of the system has been extended into a complete hospital digital ecosystem.

The first version provided a role-based structure for three types of users, namely \textbf{Patients}, \textbf{Doctors} and \textbf{Administrators}. Patients register, maintain health profiles, book specialization-based appointments from real-time slots and download prescriptions. Doctors manage their appointments and upload prescriptions. Administrators manage users, doctors, appointments and slots. \textbf{MediBot}, a chatbot built on the Google Gemini API, guides patients from their symptoms to the right specialist and books appointments in conversation.

In this phase, seven new modules have been designed, implemented and tested: (1) a \textbf{Blood Bank} with donor records, ABO/Rh-compatible issue and first-expiry-first-out stock; (2) \textbf{Bed Management} in which doctors request beds and discharges and the administrator assigns them, guarded against double allocation at the database level; (3) \textbf{Emergency Management} using the five-level Emergency Severity Index (ESI) for triage; (4) \textbf{Teleconsultation} through video rooms with a timed join window, a waiting room from which the doctor admits the patient, and signed room tokens; (5) in-app \textbf{Notifications and Alerts}; (6) \textbf{Hospital Navigation} that computes shortest and wheelchair-accessible indoor routes with Dijkstra's algorithm; and (7) \textbf{Online Payments} through Razorpay with HMAC-SHA256 signature verification and idempotent webhooks.

Security has also been strengthened. Tokens now live in httpOnly cookies with short-lived access tokens and refresh-token rotation, every state-changing request is protected by a CSRF double-submit token, all inputs are validated with Zod schemas, failed logins are rate-limited per address and per account, and uploaded medical files are stored in the database and served only to their patient and doctor. Appointments can be cancelled by any party, which frees the slot and flags paid visits for refund. The database has grown from six to twenty-one tables. The system is verified by 383 automated tests (165 backend unit, 80 backend integration against a real MySQL database, 126 frontend component and contract tests, and 12 browser end-to-end tests), and runs in a GitHub Actions pipeline. It is deployed with Docker on Render, with the frontend on Vercel and the database on managed MySQL.
